\scanpage{40}
3.1.9 Soient $\mathcal A$ et $\mathcal B$ \deleted{d'un} deux faisceaux d'anneaux de torsion sur $S$ et soit $u:\deleted{\unsure{\mathcal A}}\mathcal B\longrightarrow\deleted{\unsure{\mathcal B}}\mathcal A$ un homomorphisme. Prenons dans 3.1.8, pour $L$, le complexe de $\mathcal A,\mathcal B$-bimodules \deleted{$\mathcal B$} reduit à $\deleted{\unsure{\mathcal B}}\mathcal A$ placé en degré 0. Le foncteur
\[
 \deleted{\underline F}\ \mathrm{Mod}(S,\mathcal A)\longrightarrow\mathrm{Mod}(S,\mathcal B):
 \underline F\longmapsto\underline{\mathrm{Hom}}_{\mathcal A}(L,\underline F),
\]
est le foncteur de restriction des scalaires $u_{.}$ de $\mathcal A$ à $\mathcal B$. La flèche d'adjonction, est ici un isomorphisme de foncteurs rendant commutatif le diagramme suivant.
\diagram{\begin{tikzcd}[ampersand replacement=\&,row sep=large]
 D^+\deleted{\mathrm{Mod}}(X,f^*\mathcal A) \arrow[r,"u_{.}"]
 \& D^+\deleted{\mathrm{Mod}}(\deleted{\unsure{S}}X,f^*\mathcal B)\\
 D^+\deleted{\mathrm{Mod}}(S,\mathcal A) \arrow[r,"u_{.}"] \arrow[u,"Rf^!"']
 \& D^+\deleted{\mathrm{Mod}}(S,\mathcal B)\,. \arrow[u,"Rf^!"']
\end{tikzcd}}
On voit que le faisceau $\mathcal A$ ne joue qu'un rôle bidon dans la définition de $Rf^!$.
% slanted
\margin{\illegible{3} $\searrow$}
% archive sheet number: 39
